\documentclass[10pt,a4paper]{article}
\usepackage{fontspec,amsmath,amssymb,graphicx}
\usepackage[margin=18mm]{geometry}
\setmainfont{Times New Roman}
\setlength{\parindent}{0pt}
\setlength{\parskip}{5pt}
\allowdisplaybreaks
\sloppy
\emergencystretch=2em
\begin{document}
\section*{VD\_3\_3\_2\_03 — Lời giải bổ sung}
\subsection*{Kiến thức và trình tự giải}
\textbf{Bước 1 — Tách biến và chọn bài toán trị riêng theo biên.} Đặt $u=X(x)T(t)$ trong $u_{tt}=a^2u_{xx}$, suy $X''+\lambda X=0$, $T''+a^2\lambda T=0$. Tích phân từng phần cho $\lambda\int X^2=\int (X')^2$ khi hạng biên bằng không; do đó $\lambda\ge0$. Giải cả trường hợp $\lambda=0$ trước khi xét $\lambda=k^2>0$.

\textbf{Bước 2 — Tìm hệ hàm riêng.} Trên $(0,L)$: Dirichlet hai đầu dùng $\sin(n\pi x/L)$, $n\ge1$; Neumann hai đầu dùng $1$ và $\cos(n\pi x/L)$; Dirichlet–Neumann dùng $\sin((n+1/2)\pi x/L)$, $n\ge0$; Neumann–Dirichlet dùng cos của các tần số nửa nguyên ấy. Chuẩn bình phương của mode khác không là $L/2$, của mode hằng là $L$.

\textbf{Bước 3 — Chiếu cả hai dữ liệu.} Với mode khác không đặt


\[
A_n=\frac2L\int_0^L fX_n\,dx,\quad B_n=\frac2L\int_0^L gX_n\,dx,\quad \omega_n=ak_n.
\]


Tích phân từng khoảng nếu dữ liệu chia đoạn; dùng tích phân từng phần hoặc công thức tích–tổng khi là đa thức/lượng giác. Hệ số $B_n$ là hệ số của vận tốc, nên hệ số đứng trước $\sin\omega_nt$ phải là $B_n/\omega_n$.

\textbf{Bước 4 — Ghép nghiệm và xét mode hằng.}


\[
u=\sum_n\left(A_n\cos\omega_nt+\frac{B_n}{\omega_n}\sin\omega_nt\right)X_n.
\]


Với Neumann hai đầu, thêm $A_0+B_0t$, trong đó $A_0=L^{-1}\int f$, $B_0=L^{-1}\int g$. Với tần số bị thay đổi bởi hạng $u$ hoặc cản $u_t$, giải lại ODE của từng mode; không dùng nguyên tần số $ak_n$.

\textbf{Bước 5 — Kiểm tra biên, dữ liệu và hội tụ.} Mỗi mode thỏa biên; tại $t=0$ hai chuỗi cho khai triển của $f,g$. Nếu dữ liệu không tương thích ở góc, giải theo nghĩa năng lượng hoặc từng miền và nêu đường mất tính trơn; không đạo hàm một chuỗi yếu tùy ý. Chu kỳ được suy từ tất cả mode có hệ số khác không, kể cả mode hằng; hai tần số có tỷ số vô tỷ ngăn một chu kỳ chung.

\textbf{Chi tiết tính tích phân hệ số.} Với $k=n\pi/L$, hai nguyên hàm thường dùng là


\[
\int x\sin kx\,dx=-x\cos kx/k+\sin kx/k^2,\qquad
\int x\cos kx\,dx=x\sin kx/k+\cos kx/k^2.
\]


Với đa thức $P$, ghi cả hai lần tích phân:


\[
\int_0^L P\sin kx\,dx
=\left[-P\cos kx/k\right]_0^L+
\left[P'\sin kx/k^2\right]_0^L-
\frac1{k^2}\int_0^L P''\sin kx\,dx.
\]


Chỉ bỏ hạng biên sau khi thế $0,L$ và điều kiện thực sự của $P$. Nếu dữ liệu chia đoạn, áp dụng công thức trên từng đoạn, rồi cộng các hạng tại điểm chia. Cuối cùng nhân nghịch đảo chuẩn của hàm riêng; nếu hệ số là của vận tốc, chia thêm tần số khi ghép nghiệm. Với mẫu số bằng $0$ trong tích–tổng, quay lại tích phân gốc để tính riêng.
\subsection*{Lời giải áp dụng cho bài}
\textbf{Bước 1. Trị riêng Dirichlet–Neumann.} $X(0)=0$, $X'(L)=0$. Tích phân từng phần cho $\lambda\int X^2=\int(X')^2>0$ vì nghiệm hằng thỏa $X(0)=0$ chỉ là $0$. Với $\lambda=k^2$, $X=C\sin kx$, $\cos kL=0$, nên $k_n=(n+1/2)\pi/L$, $n\ge0$.

\textbf{Bước 2. Trực giao và hệ số.} Với hai trị riêng, trừ $X_m''X_n-X_n''X_m$ và tích phân. Hạng biên $[X_m'X_n-X_n'X_m]_0^L=0$, nên tích vô hướng bằng $0$ nếu khác trị riêng. Chuẩn là $L/2$, suy $A_n=(2/L)\int f\sin k_nx$, $B_n=(2/L)\int g\sin k_nx$.

\textbf{Bước 3. Ghép nghiệm.} $u=\sum_{n\ge0}(A_n\cos ak_nt+B_n\sin ak_nt/(ak_n))\sin k_nx$. Mỗi mode đổi dấu sau $2L/a$ và có chu kỳ $4L/a$. Tính duy nhất theo năng lượng: thông lượng ở $0$ hết vì $u_t=0$, ở $L$ hết vì $u_x=0$.

\textbf{Bước 4. Dữ liệu cụ thể.} $a=1$, $L=\pi$, $f=x$, $g=\sin^3(x/2)$. Với $k_n=n+1/2$, tích phân từng phần cho $\int_0^\pi x\sin k_nx=(-1)^n/k_n^2$, nên $A_n=2(-1)^n/(\pi k_n^2)$. Từ công thức sin lập phương, $B_0=3/4$, $B_1=-1/4$:
\[u=\sum_{n\ge0}\frac{2(-1)^n}{\pi k_n^2}\cos k_nt\sin k_nx
+\frac32\sin(t/2)\sin(x/2)-\frac16\sin(3t/2)\sin(3x/2).\]
Chu kỳ cơ bản $4\pi$. $f'(\pi)=1$ trái biên Neumann, nên nghiệm không $C^1$ qua mặt sóng từ góc phải; đây là nghiệm năng lượng/từng miền, không nghiệm cổ điển toàn bao đóng.

\textbf{Chi tiết điều kiện biên, chuẩn và tích phân dữ liệu.} Từ $X''+k^2X=0$, $X=A\cos kx+B\sin kx$. $X(0)=0$ cho $A=0$; $X'(L)=Bk\cos kL=0$ với $Bk\ne0$ cho $\cos kL=0$. Chuẩn
\[\int_0^L\sin^2k_nx\,dx=\left[x/2-\sin(2k_nx)/(4k_n)\right]_0^L=L/2.\]
Với $L=\pi$,
\[\int_0^\pi x\sin k_nx\,dx=
\left[-x\cos(k_nx)/k_n+\sin(k_nx)/k_n^2\right]_0^\pi
=(-1)^n/k_n^2.\]
Đã dùng $\cos(k_n\pi)=0$, $\sin(k_n\pi)=(-1)^n$. Khai triển $g=3\sin(x/2)/4-\sin(3x/2)/4$ rồi chiếu lên hệ sine trực giao cho $B_0=3/4,B_1=-1/4$.

\textbf{Chi tiết trực giao và dấu trị riêng.} Nhân $-X''=\lambda X$ bởi $X$, tích phân:
\[\lambda\int_0^LX^2dx=-[X'X]_0^L+\int_0^L(X')^2dx=\int_0^L(X')^2dx.\]
Hạng biên hết vì mỗi đầu có $X=0$ hoặc $X'=0$. Nếu $\lambda=0$, $\int(X')^2=0$ cho $X$ hằng; khi có một đầu Dirichlet, hằng đó bằng $0$. Với hai trị riêng,
\[(\lambda_m-\lambda_n)\int_0^LX_mX_n\,dx
=\int_0^L(-X_m''X_n+X_n''X_m)dx
=[-X_m'X_n+X_n'X_m]_0^L=0.\]
Nếu $\lambda_m\ne\lambda_n$, chia hiệu trị riêng cho tích vô hướng bằng $0$.\end{document}